\section{Approach}

\begin{figure}
\centering
\includegraphics[width=0.7\textwidth]{figs/overall.jpeg}
\vspace{-10pt}
\caption{The overview of \Name{}.} \label{frame}
\vspace*{-10pt}
\end{figure}

As shown in Figure~\ref{frame}, we propose \textbf{ConFL}, 
a hierarchy-guided framework for concurrent fault localization based on structured concurrency knowledge.
Given a bug report as the sole input, ConFL first extracts concurrency-relevant information from source code and organizes it into a Concurrent Knowledge Base (CKB). 
The CKB captures concurrency semantics at multiple abstraction levels, including components, packages, thread entry points, methods, and interaction-level behaviors.

Fault localization is then carried out as a progressive refinement process.
Starting from the bug report, an LLM incrementally queries the CKB to narrow the search space from coarse-grained components to fine-grained concurrency interactions, while keeping the exposed context bounded.
Finally, the LLM reasons over interaction-level DSL representations to rank suspicious methods and generate concurrency-aware explanations.



\subsection{Interaction DSL Generation}
\label{sec:interaction-dsl}

At the core of our approach is an \emph{interaction DSL} designed to explicitly represent concurrency context in a lightweight and comparable manner across bug reports and source code. The goal of the DSL is to capture concurrency-relevant semantics while filtering out unrelated program details, thereby enabling focused and explainable fault localization.

Recent work has demonstrated the potential of LLMs for fault localization, but a fundamental challenge remains determining what constitutes relevant context. This challenge becomes particularly severe in concurrent systems, where relevant behaviors may span multiple threads, shared resources, and asynchronous execution paths. While general fault localization often requires reasoning over deep call chains and large code regions, we observe that concurrent faults typically exhibit a more structured form of context centered around shared-variable accesses and cross-thread interactions. This motivates a representation that explicitly models concurrency-relevant interactions while suppressing unrelated program information.

\vspace{-0.5em}
\begin{table}[H]
\centering \scriptsize
\caption{Analysis Domain for Interaction DSL}
\vspace{-8pt}
\label{tab:entity-domain}
\renewcommand{\arraystretch}{1.25}
\begin{tabular}{@{}l l l@{}}
\toprule
\textbf{Entity} & \textbf{Symbol} & \textbf{Definition} \\
\midrule
method
  & $m \in \mathbb{M}$ 
  &  \\

variable
  & $v \in \mathbb{V}_{s}$ 
  &  $::= Field \cup Param$\\

lock 
  & $l \in \mathbb{L}$ 
  & $::= synchronized$ \\

thread entry point
  & $t \in \mathbb{T}_e$ 
  & $\subseteq \mathbb{M}$  \\
\bottomrule
\end{tabular}
\end{table}
\vspace{-8pt}

Based on this observation, we propose an interaction DSL that models concurrency at the granularity of shared-variable accesses across thread entry points. For each thread entry, the DSL records shared-variable accesses reachable from its downstream methods, and an interaction is constructed whenever two thread entries may access the same shared resource, whether through the same method or through different methods. Rather than attempting to reconstruct concrete schedules or happens-before relations, the DSL conservatively captures potential cross-thread interactions that may contribute to faults. Each interaction summarizes who may execute concurrently, what resources are shared, where the accesses occur, and the associated synchronization context.

The DSL is intentionally designed as a lightweight program-level abstraction rather than a predictor of bug-specific interactions. The CKB is constructed offline using conservative static analysis to capture methods reachable from thread entry points that access shared variables, which may over-approximate actual runtime behaviors. Instead of treating all extracted interactions as buggy candidates, ConFL further relies on hierarchy-guided retrieval and LLM-based reasoning to identify interactions most relevant to the bug report. This design enables scalable fault localization without requiring execution traces or reproducible thread schedules.

\Name{} targets thread-based concurrency in Java, including explicit \texttt{Thread} creation, \texttt{Executor} submissions, asynchronous callbacks, and framework-managed entry methods. Synchronization is modeled uniformly through both \texttt{synchronized} constructs and \texttt{java.util.concurrent.Lock} primitives. By representing concurrency through shared-variable accesses across thread entry points rather than through language-specific syntax, the DSL naturally supports a broad range of executor-based and asynchronous concurrency patterns commonly used in modern Java systems.

\paragraph{Concurrency Entity Extraction.}
The DSL is grounded in a small set of concurrency-relevant entities, as summarized in Table~\ref{tab:entity-domain}.
All entities correspond to statically available program constructs and can be extracted using lightweight static analysis.
Variables include both class fields and method parameters that may escape their defining scope.
Synchronization via the \texttt{synchronized} keyword is treated separately, as it appears as a method or block modifier, while other synchronization primitives (e.g., \texttt{ReentrantLock}) are modeled as shared variables.
Thread entry points are identified as methods that may initiate concurrent execution, including explicit thread creation, executor submissions, asynchronous callbacks, and framework-managed entry methods.


\begin{wrapfigure}{r}{0.35\textwidth} 
    \centering 
    \includegraphics[width=0.35\textwidth]{figs/dsl.pdf}
    \caption{Interaction DSL}\label{fig:interaction-dsl}
\end{wrapfigure}


Based on these extracted entities, we define an interaction DSL to represent potential concurrent behaviors.
Figure~\ref{fig:interaction-dsl} illustrates the structure of the Interaction DSL.
Each interaction is defined by a pair of thread entry points and a set of shared access summaries.
For a given entry pair, the DSL records the methods reachable from each entry that access shared variables, along with the corresponding synchronization context.
By organizing concurrency information at the granularity of entry pairs and access points, the DSL provides a compact yet expressive representation of concurrency-relevant context suitable for subsequent analysis and reasoning.

As shown in Algorithm~\ref{alg:interaction-extraction}, the construction of the Interaction DSL proceeds in three phases.

In the first phase, the algorithm extracts base concurrency-relevant entities from the source code, including methods, shared variables, and thread entry points.
A static call graph is then constructed to support reachability analysis.
These entities form the foundation for subsequent interaction modeling and correspond directly to the analysis domain defined earlier.

In the second phase, the algorithm performs thread-entry–centered reachability analysis.
Starting from each thread entry point, it computes the set of methods that may be executed as part of the corresponding concurrent context.
For each reachable method, the algorithm records all accesses to shared variables.
This step associates downstream shared accesses with their originating thread entries, enabling the DSL to capture cross-class and cross-call concurrency behavior.

In the final phase, the algorithm constructs interaction instances by examining pairs of thread entry points.
For each entry pair, it identifies methods that are reachable from both entries and aggregates their shared variable accesses.
An interaction is created whenever two thread entries may access shared state through one or more common methods.
Each interaction summarizes the potential concurrent behavior between the two entries by grouping the relevant methods and shared accesses into a single DSL element.

Overall, the algorithm conservatively over-approximates concurrency interactions: it captures all shared accesses that may occur under some interleaving, without committing to specific execution orders or schedules.
This design ensures soundness while keeping the interaction representation compact and suitable for scalable fault localization and LLM-based reasoning.

\vspace{-0.5em}
{\footnotesize
\begin{algorithm}[h]
\caption{Interaction DSL Extraction}
\label{alg:interaction-extraction}
\begin{algorithmic}[1]
\Require Source code $\mathcal{P}$
\Ensure Set of interactions $\mathcal{I}$

\State \textbf{// Phase 1: Extract base entities}
\State $\mathbb{M} \gets \textit{ExtractMethods}(\mathcal{P})$
\State $\mathbb{V}_s \gets \textit{ExtractVars}(\mathcal{P})$
\State $\mathbb{T}_e \gets \textit{IdentifyThreadEntries}(\mathbb{M})$
\State $\mathcal{G} \gets \textit{BuildCallGraph}(\mathbb{M})$

\State \textbf{// Phase 2: Compute reachability and accesses per thread entry}
\ForAll{$t \in \mathbb{T}_e$}
    \State $\textit{Reach}(t) \gets \textit{ComputeReachable}(\mathcal{G}, t.m)$
    \ForAll{$m \in \textit{Reach}(t)$}
        \State $\textit{Accesses}(t, m) \gets \textit{ExtractAccesses}(m, \mathbb{V}_s)$
    \EndFor
\EndFor

\State \textbf{// Phase 3: Build interactions for each thread entry pair}
\State $\mathcal{I} \gets \emptyset$
\ForAll{$(t_1, t_2) \in \mathbb{T}_e \times \mathbb{T}_e$ where $t_1 < t_2$}
    \State $M_{common} \gets \textit{Reach}(t_1) \cap \textit{Reach}(t_2)$
    \State $\textit{SharedMethods} \gets \emptyset$
    \ForAll{$m \in M_{common}$}
        \State $A_m \gets \textit{Accesses}(t_1, m) \cup \textit{Accesses}(t_2, m)$
        \If{$A_m \neq \emptyset$}
            \State $\textit{SharedMethods} \gets \textit{SharedMethods} \cup \{(m, A_m)\}$
        \EndIf
    \EndFor
    \If{$\textit{SharedMethods} \neq \emptyset$}
        \State $I \gets \textit{BuildInteraction}(t_1, t_2, \textit{SharedMethods})$
        \State $\mathcal{I} \gets \mathcal{I} \cup \{I\}$
    \EndIf
\EndFor

\State \Return $\mathcal{I}$
\end{algorithmic}
\end{algorithm}
}
\vspace{-0.5em}




\subsection{Concurrent Knowledge Base}
\label{sec:ckb}

Based on the extracted concurrency entities and interaction DSL described in the previous section, we construct a \emph{Concurrent Knowledge Base (CKB)} to organize concurrency-relevant information in a hierarchical and queryable form.
The CKB serves as the structured knowledge layer that supports hierarchical retrieval and LLM-guided fault localization.

Formally, the Concurrent Knowledge Base is defined as:
\[
\mathcal{K} = \{ \mathcal{C}, \mathcal{P}, \mathcal{E}, \mathcal{M}, \mathcal{I} \}
\]
where each element captures concurrency knowledge at a different abstraction level.

The Concurrent Knowledge Base (CKB) organizes concurrency-relevant information at multiple abstraction levels to support progressive fault localization.
At coarse granularity, component- and package-level knowledge captures high-level system structure and functionality.
Each component and package is associated with brief descriptions summarizing its responsibilities and concurrency relevance, derived from build configurations, directory structure, and developer documentation.
This information provides early-stage guidance for narrowing the search space when analyzing large codebases.

At finer granularity, the CKB records detailed concurrency semantics at the thread-entry, method, and interaction levels.
Formally, thread-entry knowledge corresponds to the set $\mathbb{T}_e$, where each $t \in \mathbb{T}_e$ represents a concurrency entry point such as a thread, executor submission, or asynchronous callback.
For each thread entry $t$, the CKB records the set of other entries $t' \in \mathbb{T}_e$ with which $t$ may interact through shared resource accesses.
Method-level knowledge is defined over the set of methods $\mathbb{M}$.
For each method $m \in \mathbb{M}$, the CKB records whether $m$ serves as a thread entry ($m \in \mathbb{T}_e$), the set of shared variables $v \in \mathbb{V}_s$ accessed by $m$, and the corresponding access types (e.g., read or write).
This information captures how concurrency behavior manifests at the code level.
Finally, interaction-level knowledge consists of interaction DSL elements $\mathcal{I}$.
Each interaction $I \in \mathcal{I}$ summarizes a potential concurrent behavior between a pair of distinct thread entries $(t_1, t_2)$, where $t_1, t_2 \in \mathbb{T}_e$.

\vspace{-0.5em}
\begin{table}[h]
\centering
\caption{Function calls for hierarchical retrieval}
\vspace{-8pt}
\label{tab:function_calls}
\scriptsize
\begin{tabular}{l|l|p{7.5cm}}
\hline
\textbf{Name} & \textbf{Arguments} & \textbf{Description} \\
\hline
list\_components & None &
Return all components with brief descriptions. \\
\hline
list\_packages\_under\_component & component\_name &
Return all packages under the specified component. \\
\hline
list\_entries\_under\_package & package\_name &
Return all thread entry points defined in the package. \\
\hline
list\_method\_under\_class & class\_name &
Return all methods under the specified class. \\
\hline
get\_entry\_summary & entry\_name &
Return concurrency summary of the thread entry and its interacting entries. \\
\hline
get\_method\_summary & method\_name &
Return concurrency-related information of the method. \\
\hline
get\_interaction\_dsl & entry\_pair &
Return interaction DSL for the given pair of thread entries. \\
\hline
get\_code\_snippet & method\_name &
Return code snippet of the specified method. \\
\hline
exit & None &
Terminate the retrieval process. \\
\hline
\end{tabular}
\end{table}
\vspace{-8pt}
\vspace{-0.5em}

\paragraph{Retrieval Functions}
\label{subsec:ckb-retrieval}

To enable hierarchical and controlled access to the Concurrent Knowledge Base, we expose a small set of retrieval functions, summarized in Table~\ref{tab:function_calls}.
Each function corresponds to a specific abstraction level in the knowledge base, ranging from components and packages to thread entries, methods, and interaction DSL elements.
By restricting retrieval to these predefined functions, the framework ensures that the LLM incrementally explores only concurrency-relevant context, preventing uncontrolled traversal of the codebase and keeping the reasoning process focused and interpretable.

\subsection{Explainable Fault Localization}
\label{sec:explainable-fl}

Given a bug report, our fault localization procedure follows a \emph{progressive refinement} strategy over the Concurrent Knowledge Base (CKB).
Rather than exposing large code fragments or entire program graphs to an LLM, we localize concurrency faults by performing hierarchical retrieval and reasoning over structured, DSL-encoded concurrency knowledge.

The hierarchical design is motivated by the structural properties of large software systems.
Well-engineered systems are typically organized with high cohesion and low coupling: components differ significantly in functionality, packages encapsulate closely related behaviors, and thread entry points define distinct concurrent execution contexts.
As a result, concurrency faults are usually confined to a small region of the system and rarely span arbitrary components or unrelated execution contexts.
This observation motivates a coarse-to-fine localization strategy that progressively narrows the search space along natural architectural boundaries.

\paragraph{Hierarchical Localization Workflow.}
Starting from the bug report, the LLM incrementally refines the localization scope by querying the CKB in the following order:
(i) selecting the most relevant component,
(ii) selecting a package within the chosen component,
(iii) localizing candidate thread entry points within the package, and
(iv) selecting a pair of interacting thread entries.
At each stage, only a small and relevant slice of the CKB is exposed to the LLM (e.g., a list of components, packages, or thread entries), ensuring that the reasoning context remains bounded and focused.

Once a thread-entry pair is selected, the framework retrieves all interaction-level records associated with the pair.
These records are represented using the interaction DSL and summarize shared resources, access actions, and synchronization context when available.
The interaction set provides a compact, concurrency-focused context that identifies the candidate methods involved in the interaction.

\paragraph{LLM-based Interaction Ranking and Explanation.}
Given the interaction DSL elements, the LLM performs interaction-level reasoning to assess which methods are most likely responsible for the fault.
To support this reasoning, the LLM may selectively query method-level summaries and, when necessary, corresponding code snippets through predefined retrieval functions.
Based on the combined interaction context and method-level information, the LLM produces a ranked list of suspicious methods.

Finally, the LLM generates an explanation grounded in explicit concurrency context.
The explanation links the ranked methods to their accessed shared variables and the corresponding thread-entry pair, describing how concurrent execution and shared-resource access may lead to the observed faulty behavior.

\begin{figure}[H] \scriptsize
\centering 
\vspace*{3pt}
\begin{tcolorbox}[colback=gray!5,colframe=black!30]

You are an expert in concurrent fault localization. Given a bug report and access to a Concurrent Knowledge Base (CKB), iteratively retrieve relevant concurrency information through hierarchical reasoning and identify the most suspicious methods related to the reported fault.

\vspace{0.3em}

\textbf{Bug Report:} \verb!{{BUG_REPORT}}!

\# Component-Level Retrieval

First, call: \textit{$list\_components()$}. Use function calls to retrieve concurrency-related knowledge from the CKB, then analyze the returned components and identify the most relevant component associated with the bug report.

\begin{verbatim}
{{function_call}}
\end{verbatim}
\textbf{Output Format:} component\_name: component
\vspace{0.3em}

\# Package-Level Retrieval

First, call: \textit{$list\_packages\_under\_component(component)$}. Use function calls to retrieve concurrency-related knowledge from the CKB, then analyze the returned packages and identify the most relevant package associated with the bug report. 

\begin{verbatim}
{{function_call}}
\end{verbatim}
\textbf{Output Format:} package\_name: package
\vspace{0.3em}

\# Thread Entry and Method-Level Retrieval

First, call: \textit{$list\_entries\_under\_package(package)$}. Use function calls to retrieve concurrency-related knowledge from the CKB, then analyze the returned thread entry points and identify the most suspicious entry pair related to the reported fault. 

Select the most suspicious thread-entry pair.

\begin{verbatim}
{{function_call}}
\end{verbatim}
\textbf{Output Format:} entry\_pair: (entry1, entry2)
\vspace{0.3em}

\# Interaction DSL Reasoning

First, retrieve the interaction DSL: $get\_interaction\_dsl(entry_pair)$. Use the DSL to analyze shared resources, access actions, synchronization context, and cross-thread interactions. If more evidence is needed, continue using function calls to inspect relevant method summaries or source code:
\begin{verbatim}
{{function_call}}
\end{verbatim}

After sufficient evidence has been collected, apply the reasoning rules to identify suspicious methods involved in the concurrent interaction. Then return a ranked list of suspicious methods and explain how the interaction may lead to the reported fault.

\textbf{Output Format:}

\begin{verbatim}
Rank 1: [Method]: ... [Shared Resource]: ... [Thread Entries]: ... [Explanation]: ...
\end{verbatim}
\end{tcolorbox}
\vspace{-0.4em}
\caption{Prompt Template} \label{template}
\vspace{-0.15in}
\end{figure}


\paragraph{Prompt Template.} 
ConFL follows the function-calling paradigm common to recent LLM-based FL approaches (e.g., AutoFL, FlexFL). Figure~\ref{template} illustrates the complete prompt template. At each hierarchical retrieval stage, the LLM is provided with the bug report and a typed set of retrieval functions (listed in Table~\ref{tab:function_calls}); it selects one function to invoke, observes the structured response, and progresses to the next stage. This constrains the LLM to operate within the CKB schema rather than generating free-form responses. At the interaction-ranking stage, each interaction DSL element is serialized as a labeled block listing the entry pair, the set of shared methods, and the per-method shared-variable accesses with their guard locks, so that the LLM can reason directly over concurrency-relevant context without traversing raw code. We use deterministic decoding (temperature = 0) and the complete prompt templates are released in the artifact.



 

